% Input in supplementary material, after \appendix. Requires theory_sections.tex.
\section{Proofs and an allocation counterexample}
\label{app:aisi_theory}

\paragraph{Bayes decision and Gaussian update.}
Conditional on $D$, the oracle term in Equation~\eqref{eq:aisi_loss} is
independent of the chosen set. Thus minimizing expected loss is equivalent
to minimizing $\sum_{i\in S}m_i$ subject to the fixed quotas. A deterministic
index rule may resolve equal means without changing the objective.
For a Gaussian linear observation $W=h^{\mathsf T}(\theta,b)+\epsilon$,
write $V=h^{\mathsf T}Ph+r>0$. Conditioning yields
\[
 P^+=P-Phh^{\mathsf T}P/V,\qquad
 m^+=m+(Ph)_1\{W-h^{\mathsf T}(m,a)\}/V.
\]
The predictive residual is $N(0,V)$, which proves
Equation~\eqref{eq:aisi_innovation}. This calculation conditions on fixed
hyperparameters; learning a shared hyperparameter from the new observation
can also update other units and invalidate the one-coordinate simplification.

\begin{proof}[Proof of Proposition~\ref{prop:aisi_kg}]
Order the other units' means as $a_{(1)}\le\cdots\le a_{(|g|-1)}$.
With $t_i=a_{(k_g)}$, the optimal posterior-mean sum in this group equals
$\sum_{j<k_g}a_{(j)}+\min(m_i,t_i)$. After the observation, replace
$m_i$ by $m_i+sZ$, where $Z\sim N(0,1)$; all other group contributions
are unchanged. The tower property cancels the expected oracle-loss term
from the before/after Bayes risk. Therefore the value is
\[
 K^{-1}\{\min(m_i,t_i)-\mathbb E\min(m_i+sZ,t_i)\}
 =K^{-1}\mathbb E(sZ-|m_i-t_i|)_+.
\]
Integrating the Gaussian tail gives Equation~\eqref{eq:aisi_kg}.
For $s>0$ and $\Delta\ge0$, differentiation gives
$\partial(K\operatorname{KG})/\partial s=\varphi(\Delta/s)>0$ and
$\partial(K\operatorname{KG})/\partial\Delta=-\Phi(-\Delta/s)<0$.
At $s=0$, the mean does not change and the value is zero. For a quota
of zero or the entire group, the feasible choice in that group is fixed;
the tower property gives zero expected improvement. Equal means cause no
problem because the optimum sum, rather than a discontinuous membership
indicator, is the value function.
\end{proof}

\begin{proof}[Proof of Proposition~\ref{prop:aisi_floor}]
The sample mean is sufficient and has variance
$\tau^2+\beta^2+r/n$, with covariance $\tau^2$ with $\theta$.
Gaussian conditioning proves Equation~\eqref{eq:aisi_level_floor}.
For two units, set $T_i=\theta_i+b_i$. Their conditional means
$m_i=\mu+\lambda(T_i-\mu)$, with
$\lambda=\tau^2/(\tau^2+\beta^2)$, are independent normals with
standard deviation $\tau^2/\sqrt{\tau^2+\beta^2}$. The minimum of two
independent normals of common mean $\mu$ and standard deviation $s$
has expectation $\mu-s/\sqrt{\pi}$. Hence
\[
 \mathbb E\min(m_1,m_2)-\mathbb E\min(\theta_1,\theta_2)
 =\frac{\tau}{\sqrt\pi}
   -\frac{\tau^2}{\sqrt{\pi(\tau^2+\beta^2)}},
\]
which is Equation~\eqref{eq:aisi_selection_floor}. Any adaptive sequence
of ordinary observations can be simulated from exact $(T_1,T_2)$ and
independent observation and policy randomness. Exact $T$ therefore
contains at least as much decision-relevant information as that sequence;
its optimal Bayes risk is a lower bound for all such policies.
\end{proof}

\paragraph{A current-posterior version of the variance floor.}
For arbitrary positive-semidefinite $P$ with entries $(v,c,q)$, exact
knowledge of $T=\theta+b$ leaves
$v-(v+c)^2/(v+q+2c)$ when $\operatorname{Var}(T)>0$.
If $\operatorname{Var}(T)=0$, $T$ is already known and its acquisition
does not change $v$; positive semidefiniteness implies $v+c=0$.
After learning $T$ exactly, $\operatorname{Cov}(\theta,T\mid T)=0$,
so further ordinary sampling has zero target-mean innovation. Under
independent units, positive remaining $v$, a nontrivial quota, and finite
noise, ideal validation still has strictly positive one-step value.
Affordability and the amount of remaining budget are separate constraints.

\begin{proof}[Proof of Proposition~\ref{prop:aisi_shared}]
Every feasible set has the same additional sum $\sum_g k_g u_g$ under
the transformation. It appears both in the selected-set sum and the
minimized oracle sum, so it cancels. All comparisons among feasible sums
are also unchanged. If ordinary means equal $\theta_i+b_g$, this same
argument applies with $u_g=b_g$ to recover the optimal feasible sets.
\end{proof}

\paragraph{Why marginal uncertainty can be irrelevant.}
For two units with $\theta_1=\theta_2=Z$, $Z\sim N(0,1)$, and one slot,
the selection regret is always zero. An exact observation of either unit
updates both posterior means to $Z$, so the true joint information value
is zero. Updating only one coordinate incorrectly would give
$\varphi(0)=1/\sqrt{2\pi}$. This shows why the independence/one-coordinate
assumption of Proposition~\ref{prop:aisi_kg} matters when offsets or targets
have shared structure.

\begin{proof}[Proof of Proposition~\ref{prop:aisi_channel}]
At the same unit the gap is the same for both actions, and
Proposition~\ref{prop:aisi_kg} is strictly increasing in $s$.
Substituting $h_A$ and $h_R$ into Equation~\eqref{eq:aisi_innovation}
and comparing squared innovations gives Equation~\eqref{eq:aisi_switch}.
For $c=0$ and $v>0$, the equal numerators $v^2$ cancel.
For the final statement, conditional on $\theta$, exact $T=\theta+b$
and $U=\theta+d$ have independent likelihoods with variances $\beta^2$
and $\delta^2$. Adding their precisions to the prior precision proves
Equation~\eqref{eq:aisi_two_biased}. Repeated independent observation
noise can reveal $T$ and $U$, but it does not remove their persistent
offsets. If both channels instead share the same offset, their exact
values supply only one such equation.
\end{proof}

\paragraph{A fixed-budget counterexample to the cost ratio.}
There is one known alternative of value zero and one target
$\theta\sim N(0,1)$, with one slot and budget two. A cheap direct
measurement of $\theta$ has noise variance one and cost one. A more
accurate direct measurement has noise variance $0.1$ and cost two.
Both have zero initial gap. Their initial knowledge gradients are
$1/\sqrt{4\pi}=0.282095$ and
$1/\sqrt{2.2\pi}=0.380377$, respectively. Their cost ratios therefore
favor the cheap action: $0.282095>0.190188$. Only another cheap action
is affordable afterward. The two cheap measurements produce posterior
mean variance $2/3$, hence total expected improvement
$\sqrt{2/3}/\sqrt{2\pi}=0.325735$. Spending the budget on the more
accurate action yields $0.380377$, a larger value. The policy can fail
even with correct Gaussian beliefs and conditionally unbiased channels;
one-step value per cost is not the multistep budget objective.
